\documentclass[12pt,a4paper]{article}

% Plantilla académica para Overleaf. Compilar con pdfLaTeX.
% Reemplazá los textos entre corchetes por los datos de tu trabajo.
\usepackage[utf8]{inputenc}
\usepackage[T1]{fontenc}
\usepackage[spanish,es-nodecimaldot]{babel}
\usepackage{lmodern}
\usepackage[a4paper,top=2.5cm,bottom=2.5cm,left=3cm,right=3cm]{geometry}
\usepackage{microtype}
\usepackage{setspace}
\usepackage{array}
\usepackage{graphicx}
\usepackage{booktabs}
\usepackage{hyperref}
\usepackage{xcolor}
\usepackage{titlesec}
\usepackage{enumitem}

\hypersetup{colorlinks=true,linkcolor=black,urlcolor=blue,citecolor=black}
\setlength{\parindent}{1.0cm}
\setlength{\parskip}{0.25em}
\onehalfspacing
\setlist{topsep=0.2em,itemsep=0.2em}
\titleformat{\section}{\normalfont\Large\bfseries}{\thesection.}{0.4em}{}
\titleformat{\subsection}{\normalfont\large\bfseries}{\thesubsection}{0.4em}{}
\titlespacing*{\section}{0pt}{1.4em}{0.4em}

% Editá estos datos para completar la portada.
\newcommand{\institucion}{Universidad de Buenos Aires}
\newcommand{\facultad}{Facultad de Ciencias Económicas}
\newcommand{\carrera}{Tecnicatura Universitaria en Gestión y Análisis de Datos en Organizaciones}
\newcommand{\tituloTrabajo}{Trabajo Práctico Grupal}
\newcommand{\subtituloTrabajo}{La influencia de la localía en el fútbol argentino}

\begin{document}

% PORTADA
\thispagestyle{empty}
\begin{center}
    {\large\bfseries \institucion}\\[0.25em]
    {\large \facultad}\\[0.25em]
    {\normalsize \carrera}
\end{center}

\vspace{1.4cm}
\noindent\rule{\textwidth}{0.7pt}

\vspace{1.5cm}
\begin{center}
    {\LARGE\bfseries \tituloTrabajo}\\[0.8em]
    {\large \subtituloTrabajo}
\end{center}

\vspace{1.4cm}
\noindent\rule{\textwidth}{0.4pt}
\vspace{0.5em}

% Reemplazá cada valor entre corchetes. Agregá o quitá filas si hace falta.
\renewcommand{\arraystretch}{1.45}
\noindent\begin{tabular}{@{}>{\bfseries}p{0.31\textwidth}p{0.63\textwidth}@{}}
Modalidad: & [Modalidad de cursada o trabajo] \\
Código: & [Código de materia / comisión] \\
Informe: & [Tipo o número de informe] \\
Presentación: & [Formato o instancia de presentación] \\
Repositorio: & https://github.com/santiagoleguizamon10/TP-labo \\
Entrega: & 6 de noviembre 2026 \\
Integrantes: & Santiago Leguizamón, Salvador Lanus, Benjamin Aranda, Francisco Pando,
Sebastián Chernomoretz \\
Docente: &  Rita Beatriz Morrone , Juan Ignacio da Torre, Joaquin Romano Abou-Eid, Manuel Fernández San Martín, Santiago lucas garcia, Tadeo Rost Kain
 \\
\end{tabular}

\vfill
\begin{center}
    {\large Período académico segundo cuatrimestre de 2026}
\end{center}
\clearpage

% ESTRUCTURA DEL INFORME. Las extensiones son orientativas; ajustá el texto
% para cumplir la consigna y reemplazá todos los marcadores entre corchetes.

% 1. Resumen: entre 5 y 8 líneas.
\section{Resumen}
[En 5--8 líneas, sintetizá el objetivo del trabajo, la metodología empleada y los principales hallazgos. Completá los hallazgos una vez realizado el análisis.]

% Extensión orientativa: media página.
\section{Introducción}
[Explicá la motivación del estudio y formulá la pregunta principal sobre la influencia de la localía en el fútbol argentino. Presentá brevemente el conjunto de datos utilizado y su fuente.]

% Extensión orientativa: media página.
\section{Datos y metodología}
[Describí el conjunto de datos, sus variables y estructura; explicá las transformaciones realizadas y las herramientas utilizadas. Indicá cómo definiste o mediste la localía y qué resultados futbolísticos analizaste.]

% Extensión orientativa: entre una y dos páginas.
\section{Resultados}
[Insertá los gráficos exportados desde Google Colab y las tablas pertinentes. Para cada elemento, explicá qué muestra y cómo responde a la pregunta principal. No olvides indicar la fuente de los datos.]

% Ejemplo para incorporar una imagen: subí el archivo a Overleaf y quitá el símbolo %.
% \begin{figure}[htbp]
%   \centering
%   \includegraphics[width=0.85\textwidth]{[nombre-del-grafico.png]}
%   \caption{[Descripción del gráfico.]}
%   \label{fig:resultado}
% \end{figure}

% Extensión orientativa: media página.
\section{Conclusiones}
[Resumí los hallazgos principales, señalá las limitaciones del análisis y proponé posibles extensiones.]

% Extensión orientativa: un cuarto de página.
\section{Reflexión metodológica}
[Contá si se utilizó inteligencia artificial, para qué tareas y qué revisiones hicieron sobre sus respuestas. Comentá los aprendizajes del proceso y las decisiones metodológicas del grupo.]

% Sección opcional. Si no van a incluir fuentes, eliminá esta sección.
\section{Bibliografía (opcional)}
[Agregá las fuentes utilizadas para contextualizar o interpretar los resultados: artículos, informes, notas o documentación del conjunto de datos. Usá el formato de cita solicitado por la cátedra.]
\end{document}
